\section{Preliminaries}

By default, all logarithms are to the base of $2$. 
%All graphs in this paper are undirected, and are allowed to have parallel edges but not self-loops.

Let $G=(V,E,\ell)$ be an edge-weighted graph, where each edge $e\in E$ has weight (or \emph{length}) $\ell_e$. 
For a vertex $v\in V$, we denote by $\deg_G(v)$ the degree of $v$ in $G$.
For each pair $S,T\subseteq V$ of disjoint subsets, we denote by $E_{G}(S,T)$ the set of edges in $G$ with one endpoint in $S$ the other endpoint in $T$.
For a pair $v,v'$ of vertices in $G$, we denote by $\dist_{G}(v,v')$ (or $\dist_{\ell}(v,v')$) the shortest-path distance between $v$ and $v'$ in $G$.
We define the \emph{diameter} of $G$ as $\diam(G)=\max_{v,v'\in V}\set{\dist_G(v,v')}$, and we define the \emph{girth} of $G$, denoted by $\gir(G)$, as the minimum weight of any cycle in $G$.
We may omit the subscript $G$ in the above notations when the graph is clear from the context.


%\paragraph{Graph expansion and a distribution of random graphs.}
%\emph{cost} of $S$ is $\vol_{G}(S)=\sum_{u\in S}\deg_{G}(u)$ and we write $\vol(G)=\vol_{G}(V)$. 
Given a graph $G$, its \emph{conductance} 
%of the cut is defined as
is defined as 
\[\Phi(G)=\min_{S\subseteq V, S \ne\emptyset, S\ne V}
\set{
\frac{|E_{G}(S,V\setminus S)|}
{\min\set{\sum_{v\in S}\deg_G(v),\sum_{v\notin S}\deg_G(v)}}
	}.\]
We say that $G$ is a \emph{$\phi$-expander} iff $\Phi(G)\ge\phi$. 

\iffalse
\paragraph{Chernoff Bound.}
We will use the following form of Chernoff bound (see e.g.\cite{dubhashi2009concentration}).
\begin{lemma}[Chernoff Bound] \label{prop:chernoff}
	Let $X_1,\dots X_n$ be independent random variables taking values in $\{0,1\}$. Let $X$ denote their sum and let $\mu=\ex{X}$ denote the sum's expected value. Then for any $\delta>0$, 
	$$
	\pr{X > (1+\delta) \mu} < \left( \frac{e^{\delta}}{(1+\delta)^{1+\delta}}\right)^{\mu}.
	$$
\end{lemma}
\fi


